\subsection{Verso uno scheduling energy-aware con consenso esplicito}

La combinazione tra (i)~l'adozione crescente di architetture edge--cloud
ibride, (ii)~i vincoli energetici stringenti nei nodi edge e
(iii)~la disponibilità di misurazioni energetiche a granularità
container-level rende oggi possibile e necessario studiare meccanismi
di esecuzione in grado di controllare dinamicamente l'energia
associata alle funzioni serverless.

Tuttavia, la riduzione del consumo energetico ottenuta tramite
approssimazione introduce un compromesso con l'accuratezza del
risultato che non può essere gestito unilateralmente dalla piattaforma.
Tale compromesso richiede un \emph{consenso esplicito} da parte
dell'utente che invoca la funzione: solo in presenza di tale consenso
lo scheduler è autorizzato a selezionare varianti approssimate.
Questa scelta progettuale garantisce che il comportamento del sistema
resti trasparente e controllabile, distinguendo nettamente tra
ottimizzazione energetica e degradazione indesiderata della qualità.
